\section{Matemáticas formales: objeto de búsqueda, entorno y verificador}

Al pasar de la traducción a theorem proving, la tarea cambia de raíz. La generación de lenguaje natural puede aceptar varias respuestas aproximadamente correctas, mientras que una demostración formal exige que cada paso cumpla las reglas de tipos y de lógica del proof assistant. El modelo ya no solo predice tokens en el espacio del texto, sino que busca dentro de un entorno \textbf{con estado, ejecutable y capaz de rechazar acciones erróneas}.

\subsection{Por qué proof search no es un árbol de búsqueda común como el de los juegos de tablero}

El \term{formal prover} (probador formal) guarda el estado actual de la demostración y verifica las tactics; el \term{proof state} es el conjunto de objetivos aún no resueltos; una \term{tactic} es una transformación válida que se aplica a un objetivo; un \term{subgoal} es un objetivo nuevo producido por una tactic. Si se usa inducción matemática sobre una proposición universal, un objetivo produce a la vez dos subobjetivos, el caso base y el paso inductivo, y el objetivo padre solo se considera completo cuando ambos están completos.

El proof state puede escribirse como
\[
\mathcal{G}=\{g_1,g_2,\ldots,g_k\},
\qquad
a(g_i)\rightarrow\{g_{i1},g_{i2},\ldots,g_{im}\}.
\]
Aquí la acción $a$ no produce un único estado sucesor, sino que puede producir un grupo de subobjetivos que deben resolverse en conjunto. Por eso la estructura de búsqueda se parece más a un hypergraph; un hyperedge conecta el objetivo padre con varios subobjetivos necesarios.

\lecturefigure{03-hypertree-search.png}{Una tactic puede generar varios subgoals que deben cerrarse simultáneamente.}{Subtítulos locales 05:00--07:00; Lample et al., \emph{HyperTree Proof Search}, 2022.}

\begin{knowledgebox}{Lectura de la figura: AND y OR deben coexistir}
El modelo puede proponer varias tactics para un mismo objetivo; esa es la elección OR. En cuanto se elige induction, el caso base y el paso inductivo forman a su vez una restricción AND. La intuición de MCTS común, que solo contempla que «una acción produce un estado nuevo», no basta: el algoritmo de búsqueda debe acumular la probabilidad conjunta de éxito de un grupo de objetivos y propagar los fallos de verificación de vuelta a la política de selección.
\end{knowledgebox}

\subsection{Informal-to-formal: que el borrador reduzca el espacio de búsqueda}

Internet contiene una gran cantidad de matemáticas en lenguaje natural o en \LaTeX{}, pero muy poco corpus formalizado. El puente central que se presentó en clase es este: primero, un LLM de propósito general genera una informal proof; después, el teorema y ese borrador de demostración se entregan al formal model para ayudarlo a elegir tactics. Este diseño no rebaja el criterio final de corrección, porque el proof assistant sigue verificando paso a paso la demostración completa; solo aporta al buscador una heurística con mucha información.

\lecturefigure{04-informal-to-formal.png}{La demostración informal aporta la ruta; el probador formal conserva la decisión final.}{Subtítulos locales 07:00--09:10; HyperTree Proof Search y el trabajo informal-to-formal descrito oralmente en clase.}

\begin{importantbox}{Separación de responsabilidades entre generador y verificador}
El generador busca cobertura y capacidad de exploración, y puede proponer rutas incompletas o erróneas; el verificador debe comprobar las reglas de forma determinista. La confiabilidad del sistema proviene de «proponer con holgura, aceptar con rigor», no de exigir que el modelo generativo acierte por completo al primer intento. Este patrón aparece después también en la ejecución de código, los verifier matemáticos, los tool-use agent y el reasoning RL.
\end{importantbox}

\begin{warningbox}{Una informal proof no es una demostración que pueda entregarse}
Un borrador en lenguaje natural puede saltarse pasos, usar lemas no declarados o suponer implícitamente condiciones adicionales. Su valor para la formal search es comprimir el espacio de candidatos, no sustituir al proof assistant. Si un sistema solo muestra una cadena en lenguaje natural sin verificación ejecutable, no puede afirmar que ofrece una garantía formal.
\end{warningbox}

\teachervoice{Lample enfatiza que, en ese momento, la comunidad de matemáticas formales casi no tenía API preparadas para el aprendizaje automático, y el equipo tuvo que construir por su cuenta la interfaz entre el modelo y el prover. Este detalle muestra que el cuello de botella de una nueva línea de investigación muchas veces no es «la falta de un modelo más grande», sino la falta de un entorno capaz de reiniciar, ejecutar, verificar y registrar estados por lotes.}

\subsection{Por qué este proyecto generó la necesidad de un LLM propio}

En los primeros experimentos se podía llamar a la Codex API externa para generar informal proofs; pero si cada expansión de un nodo de búsqueda requiere una llamada remota, el costo, la latencia, las cuotas y la reproducibilidad limitan una búsqueda a escala de MCTS. Tener un modelo propio significa poder controlar el procesamiento por lotes, la temperatura de muestreo, los rollout en paralelo, la caché y las versiones, con lo que «llamar al modelo de vez en cuando» se convierte en «el modelo está dentro del ciclo interno de la búsqueda». Esa es precisamente la motivación de sistemas para pasar de la investigación aplicada al entrenamiento de modelos fundacionales.

\subsection{Resumen de la sección}

La demostración formal coloca al LLM dentro de un entorno con estado: el modelo propone tactics, el proof assistant las verifica y el algoritmo de búsqueda maneja la estructura AND/OR. El valor de informal-to-formal no está en un lenguaje más elegante, sino en usar un corpus abundante para orientar una búsqueda formal que dispone de pocos datos.

